\documentclass[10pt,a4paper]{amsart}
\usepackage{amsmath,amssymb,amsthm,graphicx}
% Retained preamble for the Stacks Project chapter morphisms.tex (Morphisms of Schemes).
% The chapter file pinned-original-source.tex carries only \input{preamble} on its first line, so no
% preamble text lives inside the chapter itself.  The Stacks Project uses one single shared preamble for
% every chapter; the snapshot stages it verbatim in sources/stacks-chow/preamble.tex.  The block below is
% copied mechanically, line for line, from that file's lines 14-18 (the Xy-pic setup) and lines 154-217
% (the native theorem declarations together with the mathematical macros).  Nothing was added, removed or
% redefined.  The wrapper class and the amsmath/amssymb/amsthm/graphicx base packages are supplied by the
% assembler, so the class/\externaldocument part of the shared preamble is not repeated here.
% For commutative diagrams we use Xy-pic
\usepackage[all]{xy}

% We use 2cell for 2-commutative diagrams.
\xyoption{2cell}
% Theorem environments.
%
\theoremstyle{plain}
\newtheorem{theorem}[subsection]{Theorem}
\newtheorem{proposition}[subsection]{Proposition}
\newtheorem{lemma}[subsection]{Lemma}

\theoremstyle{definition}
\newtheorem{definition}[subsection]{Definition}
\newtheorem{example}[subsection]{Example}
\newtheorem{exercise}[subsection]{Exercise}
\newtheorem{situation}[subsection]{Situation}

\theoremstyle{remark}
\newtheorem{remark}[subsection]{Remark}
\newtheorem{remarks}[subsection]{Remarks}

\numberwithin{equation}{subsection}

% Macros
%
\def\lim{\mathop{\mathrm{lim}}\nolimits}
\def\colim{\mathop{\mathrm{colim}}\nolimits}
\def\Spec{\mathop{\mathrm{Spec}}}
\def\Hom{\mathop{\mathrm{Hom}}\nolimits}
\def\Ext{\mathop{\mathrm{Ext}}\nolimits}
\def\SheafHom{\mathop{\mathcal{H}\!\mathit{om}}\nolimits}
\def\SheafExt{\mathop{\mathcal{E}\!\mathit{xt}}\nolimits}
\def\Sch{\mathit{Sch}}
\def\Mor{\mathop{\mathrm{Mor}}\nolimits}
\def\Ob{\mathop{\mathrm{Ob}}\nolimits}
\def\Sh{\mathop{\mathit{Sh}}\nolimits}
\def\NL{\mathop{N\!L}\nolimits}
\def\CH{\mathop{\mathrm{CH}}\nolimits}
\def\proetale{{pro\text{-}\acute{e}tale}}
\def\etale{{\acute{e}tale}}
\def\QCoh{\mathit{QCoh}}
\def\Ker{\mathop{\mathrm{Ker}}}
\def\Im{\mathop{\mathrm{Im}}}
\def\Coker{\mathop{\mathrm{Coker}}}
\def\Coim{\mathop{\mathrm{Coim}}}

% Boxtimes
%
\DeclareMathSymbol{\boxtimes}{\mathbin}{AMSa}{"02}

%
% Macros for moduli stacks/spaces
%
\def\QCohstack{\mathcal{QC}\!\mathit{oh}}
\def\Cohstack{\mathcal{C}\!\mathit{oh}}
\def\Spacesstack{\mathcal{S}\!\mathit{paces}}
\def\Quotfunctor{\mathrm{Quot}}
\def\Hilbfunctor{\mathrm{Hilb}}
\def\Curvesstack{\mathcal{C}\!\mathit{urves}}
\def\Polarizedstack{\mathcal{P}\!\mathit{olarized}}
\def\Complexesstack{\mathcal{C}\!\mathit{omplexes}}
% \Pic is the operator that assigns to X its picard group, usage \Pic(X)
% \Picardstack_{X/B} denotes the Picard stack of X over B
% \Picardfunctor_{X/B} denotes the Picard functor of X over B
\def\Pic{\mathop{\mathrm{Pic}}\nolimits}
\def\Picardstack{\mathcal{P}\!\mathit{ic}}
\def\Picardfunctor{\mathrm{Pic}}
\def\Deformationcategory{\mathcal{D}\!\mathit{ef}}

% --- appended 2026-10-12 (Chief): the project own definitions of the reference/slogan/history
% display environments, copied mechanically from the shared Stacks preamble lines 8-12
% (\usepackage{verbatim} plus the three \newenvironment lines, which the chapter preamble keeps
% verbatim from the project preamble). Nothing added, removed or redefined: these environments
% carry no mathematical content and the project itself typesets them as comments. Without these
% three lines a counted statement whose own text uses them fails with Environment slogan undefined.
\usepackage{verbatim}
\newenvironment{reference}{\comment}{\endcomment}
\newenvironment{slogan}{\comment}{\endcomment}
\newenvironment{history}{\comment}{\endcomment}
\usepackage{cleveref}
\crefname{originallemma}{original lemma}{original lemmas}
\Crefname{originallemma}{Original lemma}{Original lemmas}
\makeatletter
\newlabel{algebra-lemma-integral-closure-localize}{{0307}{source}{}{sourcepointer-1}{}}
\newlabel{algebra-lemma-integral-closure-localize@cref}{{[originallemma][0][]0307}{[1][0][]source}{}{}{}}
\makeatother
\crefname{originallemma}{original lemma}{original lemmas}
\Crefname{originallemma}{Original lemma}{Original lemmas}
\makeatletter
\newlabel{algebra-lemma-total-ring-fractions-no-embedded-points}{{02LX}{source}{}{sourcepointer-2}{}}
\newlabel{algebra-lemma-total-ring-fractions-no-embedded-points@cref}{{[originallemma][0][]02LX}{[1][0][]source}{}{}{}}
\makeatother
\crefname{originallemma}{original lemma}{original lemmas}
\Crefname{originallemma}{Original lemma}{Original lemmas}
\makeatletter
\newlabel{lemma-description-normalization}{{035P}{source}{}{sourcepointer-3}{}}
\newlabel{lemma-description-normalization@cref}{{[originallemma][0][]035P}{[1][0][]source}{}{}{}}
\makeatother
\crefname{originallemma}{original lemma}{original lemmas}
\Crefname{originallemma}{Original lemma}{Original lemmas}
\makeatletter
\newlabel{lemma-integral-closure}{{035F}{source}{}{sourcepointer-4}{}}
\newlabel{lemma-integral-closure@cref}{{[originallemma][0][]035F}{[1][0][]source}{}{}{}}
\makeatother
\title{The stalk of the normalization at a point is the integral closure in the product of residue fields of the minimal primes of the local ring}
\author{The Stacks Project Authors}
\date{stacks-project:morphisms.tex:local-snapshot-20261012; GNU FDL 1.2 or later}
\begin{document}
\maketitle
\noindent\textit{Editorial note.} All delivered mathematical text is the author's own, copied from the pinned Stacks Project chapter file pinned-original-source.tex; no statement or proof step was rewritten, reordered or removed.\par\medskip
\noindent\textit{Editorial note.} Following the multi-claim convention for this source, source\_url is the source-level URL https://stacks.math.columbia.edu/ and the source label lemma-stalk-normalization is carried in claim\_key and here rather than in a per-claim tag URL.\par\medskip
\noindent\textit{Editorial note.} The preamble is the chapter's shared Stacks Project preamble, retained as sources/stacks-morphisms/retained-preamble.tex: a line-for-line excerpt of the shared Stacks preamble staged with this snapshot. No macro was invented or redefined, and this card adds no preamble text of its own.\par\medskip
\noindent\textit{Editorial note.} Statement span 14465-14479 and proof span 14481-14501 come from the halfA pool card operations/pool-cards/stacks-morphisms-halfA.json and were re-read in the pinned chapter file before this card was written.\par\medskip
\noindent\textit{Editorial note.} The proof offers two alternative routes to the first identification (the construction of the normalization, or the description of the normalization with localization of integral closures); both are preserved as the author wrote them.\par\medskip
\noindent\textit{Editorial note.} 4 cross-references occur inside the counted spans; they are listed in external\_references as navigation pointers only.\par\medskip
\noindent\textit{Allowed standard prerequisites (editorial).}\begin{itemize}
\item The construction of the normalization and the identification of its pushforward of the structure sheaf by integral closures, cited from the same chapter. Cited at proof lines 14481-14485. Gives the identification of the stalk of the pushforward with the integral closure in a product of residue fields.
\item An alternative description of the normalization, and the compatibility of integral closures with localization, cited from the same chapter and the algebra chapter. Cited at proof lines 14485-14488. Supplies a second route to the same identification, so that the two descriptions are seen to agree.
\item For a reduced ring with finitely many minimal primes the product of the residue fields at those primes is the total ring of fractions. Cited at proof lines 14492-14495. Identifies the product appearing in the second description with the total ring of fractions of the reduced local ring.
\end{itemize}\medskip
\ifdefined\hypertarget\hypertarget{sourcepointer-1}{}\fi
\noindent\textit{Source reference.} Original lemma 0307: focused-seeds/AGP/sources/stacks-snapshot-20261011/tags.txt, source lines 3938--3938. Same-chapter reference used in the retained author text at source line(s) 13874, 14339, 14493. It is supplied only as a navigation pointer so that the retained text compiles; the counted item delivers the author's own statement and the author's complete proof, and no part of the referenced item is copied in or substituted for the counted argument.\par\medskip
\ifdefined\hypertarget\hypertarget{sourcepointer-2}{}\fi
\noindent\textit{Source reference.} Original lemma 02LX: focused-seeds/AGP/sources/stacks-snapshot-20261011/tags.txt, source lines 3424--3424. Same-chapter reference used in the retained author text at source line(s) 14451, 14498. It is supplied only as a navigation pointer so that the retained text compiles; the counted item delivers the author's own statement and the author's complete proof, and no part of the referenced item is copied in or substituted for the counted argument.\par\medskip
\ifdefined\hypertarget\hypertarget{sourcepointer-3}{}\fi
\noindent\textit{Source reference.} Original lemma 035P: focused-seeds/AGP/sources/stacks-snapshot-20261011/tags.txt, source lines 4136--4136. Same-chapter reference used in the retained author text at source line(s) 14491, 14535, 14570, 14599, 14678, 14700, 15083. It is supplied only as a navigation pointer so that the retained text compiles; the counted item delivers the author's own statement and the author's complete proof, and no part of the referenced item is copied in or substituted for the counted argument.\par\medskip
\ifdefined\hypertarget\hypertarget{sourcepointer-4}{}\fi
\noindent\textit{Source reference.} Original lemma 035F: focused-seeds/AGP/sources/stacks-snapshot-20261011/tags.txt, source lines 4126--4126. Same-chapter reference used in the retained author text at source line(s) 13884, 13969, 14199, 14489. It is supplied only as a navigation pointer so that the retained text compiles; the counted item delivers the author's own statement and the author's complete proof, and no part of the referenced item is copied in or substituted for the counted argument.\par\medskip
\begin{lemma}
\label{lemma-stalk-normalization}
Let $X$ be a scheme such that every quasi-compact open has a finite
number of irreducible components. Let $\nu : X^\nu \to X$
be the normalization of $X$. Let $x \in X$. Then the following
are canonically isomorphic as $\mathcal{O}_{X, x}$-algebras
\begin{enumerate}
\item the stalk $(\nu_*\mathcal{O}_{X^\nu})_x$,
\item the integral closure of $\mathcal{O}_{X, x}$ in
the total ring of fractions of $(\mathcal{O}_{X, x})_{red}$,
\item the integral closure of $\mathcal{O}_{X, x}$ in the product
of the residue fields of the minimal primes of $\mathcal{O}_{X, x}$
(and there are finitely many of these).
\end{enumerate}
\end{lemma}

\begin{proof}
After replacing $X$ by an affine open neighbourhood
of $x$ we may assume that $X$ has finitely many irreducible
components and that $x$ is contained in each of them.
Then the stalk $(\nu_*\mathcal{O}_{X^\nu})_x$ is the
integral closure of $A = \mathcal{O}_{X, x}$ in the product $L$
of the residue fields of the minimal primes of $A$.
This follows from the construction of the normalization and
Lemma \ref{lemma-integral-closure}.
Alternatively, you can use Lemma
\ref{lemma-description-normalization}
and the fact that normalization commutes with localization
(Algebra, Lemma \ref{algebra-lemma-integral-closure-localize}).
Since $A_{red}$ has finitely many minimal primes
(because these correspond exactly to the generic points
of the irreducible components of $X$ passing through $x$)
we see that $L$ is the total ring of fractions of $A_{red}$
(Algebra, Lemma \ref{algebra-lemma-total-ring-fractions-no-embedded-points}).
Thus our ring is also the integral closure of $A$ in the total
ring of fractions of $A_{red}$.
\end{proof}

\end{document}
